\section{Compara\c{c}\~ao com o VOF}
\label{sec:comparacao}

\subsection{Por que a compara\c{c}\~ao \'e leg\'itima}

Comparar dois m\'etodos s\'o informa se eles resolverem o mesmo problema. O
HiGFlow j\'a possu\'ia um caso VOF que, examinado, \emph{\'e} o teste de gota
est\'atica de Laplace, sem nenhuma adapta\c{c}\~ao: gota circular de raio
$R=1/6$ centrada em $(0{,}5;\,0{,}5)$, dom\'inio $[0,1]^2$, duas fases
newtonianas, tens\~ao superficial ligada e velocidade inicial nula.

Trocando em mi\'udos, os dois lados coincidem em tudo que importa:

\begin{center}
\begin{tabular}{lll}
\toprule
 & front-tracking & VOF \\
\midrule
dom\'inio            & \multicolumn{2}{c}{$[0,1]^2$, malha $61\times61$, $h=1/61$} \\
gota                 & \multicolumn{2}{c}{c\'irculo $R=1/6$ em $(0{,}5;\,0{,}5)$} \\
c\'elulas por raio   & \multicolumn{2}{c}{$10{,}2$} \\
propriedades         & \multicolumn{2}{c}{$\rho=\mu=1$ nas duas fases} \\
$\sigma$ efetivo     & \multicolumn{2}{c}{$1$} \\
$Re$, $\Delta t$, passos & \multicolumn{2}{c}{$1$; $10^{-3}$; $101$} \\
salto esperado       & \multicolumn{2}{c}{$\Delta p = \sigma/R = 6$} \\
\midrule
interface            & polilinha de 64 marcadores & fra\c{c}\~ao volum\'etrica \\
tens\~ao superficial & $\sigma\kappa\mathbf{n}$ espalhada (Roma) & CSF~\cite{brackbill1992} \\
\bottomrule
\end{tabular}
\end{center}

O $\sigma$ efetivo merece uma palavra, porque n\~ao \'e \'obvio que coincida. No
caminho multif\'asico, a for\c{c}a interfacial entra na equa\c{c}\~ao de momento
como $\mathbf{IF}/(We\,\rho)$ com $We = Ca\cdot Re$. Com $Ca=Re=1$ e $\rho=1$,
resulta $We=1$ e portanto $\sigma_{\text{ef}}=1$ --- exatamente o $\sigma$ que o
front-tracking usa explicitamente. A coincid\^encia \'e feliz e foi verificada no
c\'odigo, n\~ao suposta.

Al\'em disso: mesmo solver, mesmo passo no tempo, mesmas sondas de press\~ao nos
mesmos pontos, e as correntes par\'asitas lidas do \emph{mesmo} bloco de
impress\~ao do solver nos dois casos.

\paragraph{Um confundidor que precisou ser removido.} O caso VOF, como
distribu\'ido, imp\~oe na parede superior uma velocidade que \emph{cresce no
tempo}, $8\,(1+\tanh(8t-4))\,x^{2}(1-x)^{2}$. Com ela ligada, a velocidade medida
n\~ao \'e corrente par\'asita: \'e par\'asita mais escoamento for\c{c}ado, e
compar\'a-la com a gota est\'atica do front-tracking mediria coisas diferentes.
O for\c{c}amento foi zerado para a compara\c{c}\~ao, atr\'as de um seletor que
preserva o comportamento padr\~ao do exemplo.

\subsection{Resultados}

\pgfplotstabletypeset[
  string type,
  columns/metrica/.style={column name={m\'etrica}},
  columns/ft/.style={column name={FT espalhado}},
  columns/ftbal/.style={column name={FT balanceado}},
  columns/vof/.style={column name={VOF}},
  columns/melhor/.style={column name={melhor}},
  every head row/.style={before row=\toprule, after row=\midrule},
  every last row/.style={after row=\bottomrule},
]{\dat{comparacao-laplace.dat}}

Em valores brutos:

\pgfplotstabletypeset[
  string type,
  columns/grandeza/.style={column name={grandeza}},
  columns/ft/.style={column name={FT espalhado}},
  columns/ftbal/.style={column name={FT balanceado}},
  columns/vof/.style={column name={VOF}},
  columns/exato/.style={column name={exato}},
  every head row/.style={before row=\toprule, after row=\midrule},
  every last row/.style={after row=\bottomrule},
]{\dat{comparacao-bruto.dat}}

\subsection{Leitura dos tr\^es resultados}

\paragraph{Salto de press\~ao: o front-tracking ganha, seis vezes.}
$\Delta p = 5{,}9954$ contra $6{,}0284$, para o valor exato $6$. \'E o resultado
esperado, e a raz\~ao \'e direta: a curvatura do front-tracking sai da geometria
expl\'icita da frente e \'e exata para um c\'irculo (Se\c{c}\~ao~\ref{sec:verificacao}),
enquanto a do VOF vem de derivadas de um campo de fra\c{c}\~ao suavizado --- o
ponto historicamente fr\'agil do m\'etodo.

\paragraph{Conserva\c{c}\~ao de massa: o VOF ganha, por nove ordens de grandeza.}
A deriva no tempo do VOF \'e \num{2.9e-15}, isto \'e, zero de m\'aquina: ele
advecta um escalar conservado, e conservar \'e propriedade de constru\c{c}\~ao.
O front-tracking deriva \num{1.3e-5} em 100 passos. Tamb\'em esperado, e \'e o
outro lado da mesma moeda.

Aqui vale registrar um erro de medi\c{c}\~ao cometido e corrigido no caminho:
a primeira compara\c{c}\~ao punha os dois lados em \num{1.3e-5} e parecia empate.
Era artefato de medir coisas diferentes --- para o VOF, desvio contra a \'area
anal\'itica; para o front-tracking, deriva contra a \'area inicial. Separadas as
duas quantidades, a \num{1.3e-5} do VOF revelou-se erro de \emph{inicializa\c{c}\~ao}
(a discretiza\c{c}\~ao do c\'irculo em $t=0$) e a deriva verdadeira, zero de
m\'aquina.

\paragraph{Correntes par\'asitas: na forma espalhada o VOF ganha, dez vezes ---
e isso contraria a expectativa.} \'E o resultado que motivou a
Se\c{c}\~ao~\ref{sec:balanco}, onde a causa \'e identificada e removida: a
coluna ``FT balanceado'' das tabelas acima j\'a traz o desfecho, e vem de l\'a.
Na forma espalhada, as par\'asitas \emph{estacionam} num patamar; na balanceada,
\emph{decaem} at\'e o piso de impress\~ao do solver --- a diferen\'ca entre um
for\c{c}amento esp\'urio persistente e a aus\^encia dele.

\begin{center}
\begin{tikzpicture}
\begin{semilogyaxis}[
  width=12cm, height=6.4cm,
  xlabel={passo}, ylabel={$\max\lvert u\rvert$},
  legend pos=north east, grid=major,
  legend cell align=left,
]
\addplot table[x=passo, y=ft]    {\dat{parasitas-serie.dat}};
\addlegendentry{FT espalhado}
\addplot table[x=passo, y=vof]   {\dat{parasitas-serie.dat}};
\addlegendentry{VOF}
\addplot table[x=passo, y=ftbal] {\dat{parasitas-serie.dat}};
\addlegendentry{FT balanceado}
\end{semilogyaxis}
\end{tikzpicture}
\end{center}

\subsection{Erro de forma na inicializa\c{c}\~ao}

A linha da inicializa\c{c}\~ao merece ressalva. O d\'eficit de \num{1.6e-3} do
front-tracking \'e o d\'eficit de \'area de um pol\'igono de 64 lados inscrito no
c\'irculo --- artefato de \emph{como se inicializa}, n\~ao propriedade do
m\'etodo. Ele melhora com mais marcadores, e desapareceria ao inscrever o
pol\'igono com a \'area correta. O VOF, que integra a fra\c{c}\~ao com
subdivis\~ao $16\times16$ por c\'elula, parte muito mais pr\'oximo. A compara\c{c}\~ao
nesta linha \'e entre escolhas de implementa\c{c}\~ao, e n\~ao entre m\'etodos.
